% =============================================================================
% From Verified Hypotheses to Model Mechanisms — NAACL 2027 SHORT PAPER
% Template: official ACL style (acl.sty, acl_natbib.bst), as in the Overleaf
% "Association for Computational Linguistics ACL Conference" template.
%
% Layout: main.tex holds ONLY the preamble, title/authors, and the abstract.
% Every section lives in sections/*.tex and is pulled in via \input.
%
% Short-paper budget: 4 pages main text (Intro..Conclusion). Limitations,
% Ethics, references, and appendices do NOT count toward the limit.
% Red [TODO: ...] marks = pending experiments or material to port from the
% workshop version's source; flip \drafttrue -> \draftfalse to hide them.
% =============================================================================
\documentclass[11pt]{article}

% Change "review" to "final" for camera-ready, "preprint" for arXiv.
\usepackage[review]{acl}

% Standard ACL template packages
\usepackage{times}
\usepackage{latexsym}
\usepackage[T1]{fontenc}
\usepackage[utf8]{inputenc}
\usepackage{microtype}
\usepackage{inconsolata}
\usepackage{graphicx}

% Paper-specific packages
\usepackage{amsmath,amssymb}
\usepackage{booktabs}
\usepackage{enumitem}
\usepackage{xcolor}

% Draft machinery: red TODO marks for pending results / material to port.
\newif\ifdraft\drafttrue
\newcommand{\TODO}[1]{\ifdraft\textcolor{red!70!black}{[TODO: #1]}\fi}

\title{From Verified Hypotheses to Model Mechanisms:\\
Grounding Interpretability-Based Discovery}

% Anonymized automatically by the [review] option; fill in for camera-ready.
\author{First Author \\
  Affiliation \\
  \texttt{email@domain} \\\And
  Second Author \\
  Affiliation \\
  \texttt{email@domain} \\}

\begin{document}
\maketitle

\begin{abstract}
Causal interpretability methods such as activation patching and distributed
alignment search validate hypotheses about model mechanisms, but the
hypotheses themselves are usually supplied by researcher intuition. We study
a setting where hypothesis generation can be made independent and
verifiable. In linguistic Rosetta Stone puzzles, a model must infer an
unknown grammar from translation pairs and apply it to new sentences; an
Answer Set Programming solver derives and \emph{certifies} the minimal
grammatical rules each puzzle requires, and each certified rule becomes a
causal hypothesis about the model's internals---selected by the task, not by
the analyst. Using standard interventions on \texttt{Llama-3.1-8B} across 29
certified rules from four questions in three languages, we find that
morphosyntactic rules localize in early layers while semantic lemma
selection resolves near the output; that rank-$\le$8 linear subspaces
suffice for causal control of individual rules; and that interventions on
independent rules largely commute, consistent with additive composition.
Linear probes recover every rule at the final token, yet interventions there
have near-zero causal effect---a caution for last-token readout pipelines.
\TODO{add one sentence of hypothesis-generation statistics (certification
rates, repair iterations) once the extraction experiments complete}
\end{abstract}

\section{Introduction}
\label{sec:intro}

Mechanistic interpretability has a mature toolbox for \emph{validating}
hypotheses about model internals: activation patching quantifies where
information is causally used \citep{vig2020,meng2022,zhang2024}, and
distributed alignment search (DAS) tests whether a proposed causal variable
is implemented in a linear subspace \citep{geiger2024das,wu2023boundless}.
What these methods do not provide is the hypothesis itself. In practice the
analyst chooses a concept---gender, sentiment, a high-level algorithm---and
constructs contrastive data by hand, with no independent guarantee that the
chosen concept is the one the task actually requires
\citep{hase2023,lanham2023}. Validation is rigorous; generation is
intuition.

We study a setting where generation can be made independent and verifiable.
Rosetta Stone puzzles from the UK Linguistics Olympiad (UKLO) present a
small set of translation pairs in an unfamiliar language---for example,
Gilbertese \emph{E nooria te uaa te aomata} `The man saw the fruit'---and
ask for translations of held-out sentences \citep{bozhanov2013,lingoly}.
The puzzles are self-contained: every rule needed to answer is inferable
from the given pairs. This yields built-in verification. We use an Answer
Set Programming (ASP) solver \citep{gelfond1988,gebser2014} to derive a
minimal set of grammatical rules that provably regenerates every context
pair and uniquely translates the target question. Each \emph{certified}
rule (an object-agreement suffix, a tense particle, a noun-class concord)
is a discrete, machine-verified hypothesis about what a model solving the
puzzle must compute---selected by the task structure, not by the analyst,
and converted automatically into counterfactual minimal pairs for causal
testing.

We apply this framework to \texttt{Llama-3.1-8B} on four questions from
three typologically distinct languages (Gilbertese, Ndebele, Tshiluba),
yielding 29 certified rules. Our contributions:
\begin{enumerate}[leftmargin=*,itemsep=0pt,topsep=2pt]
  \item A framework for \textbf{certified hypothesis generation}: causal
  variables derived from task structure and verified by a solver,
  decoupling hypothesis generation from both researcher intuition and the
  model under study (\S\ref{sec:generation}).
  \item A quantitative characterization of the generation pipeline itself:
  certification rates, repair-loop dynamics, and uniqueness of the
  certified grammars. \TODO{pending extraction experiments; see
  \S\ref{sec:results-generation}}
  \item \textbf{Causal validation} using standard interventions
  (\S\ref{sec:validation}): rule implementations stratify by
  depth---morphosyntax early (L0--L4), lemma selection near the output
  (L20--L31); individual rules admit causally sufficient linear subspaces
  of rank $\le 8$; interventions on independent rules commute (100\% of 16
  rule pairs); and a dissociation---every rule is linearly decodable at the
  final token, where interventions nonetheless have near-zero causal
  effect (\S\ref{sec:results-validation}).
\end{enumerate}

We use only off-the-shelf validation machinery and claim no novelty there;
the contribution is the provenance and verifiability of the hypotheses, and
the resulting evidence about how in-context grammatical rules are
implemented.

\section{Related Work}
\label{sec:related}

\paragraph{Linguistic puzzles.} Olympiad puzzles probe rule induction from
small evidence sets \citep{lingoly,modeling,unveiling,majmudar2025};
obfuscation studies show part of measured performance is memorization
\citep{lingolytoo}. These works characterize behavior; we test whether the
inferred rules correspond to internal computation.

\paragraph{Program induction and verified reasoning.} Solvers have been
used to synthesize linguistic theories \citep{ellis2022} and to verify
LLM-proposed hypotheses or reasoning
\citep{wang2023hypothesis,parab2025,lyu2023,pan2023,olausson2023,ye2023}.
We repurpose verification: the certified program is not a way to produce an
answer but the specification of interpretability hypotheses.

\paragraph{Causal interpretability.} Probing establishes decodability but
not causal use \citep{conneau2018,hewitt2019,elazar2021,belinkov2022};
activation patching \citep{vig2020,meng2022,zhang2024} and DAS
\citep{geiger2021,geiger2024das,wu2023boundless,geiger2025} localize causal
variables, with path patching extending to circuits
\citep{wang2022ioi,goldowskydill2023,conmy2023}. The low-rank orthonormal
subspace parameterization we use is standard practice in this line, as
implemented in pyvene \citep{wu2024pyvene} and used by ReFT
\citep{wu2024reft} and AxBench \citep{wu2025axbench}. All of these methods
presuppose a known target variable; our targets come from task
certificates. Composition of features and interventions has been studied
via steering and feature geometry
\citep{turner2023,rimsky2024,merullo2024,bricken2023,huben2024,park2024};
certified rules let us test composition against operations the task
provably requires.

\section{Certified Hypothesis Generation}
\label{sec:generation}

\paragraph{Grammar schema.} A grammar is a triple
$G = (\mathcal{L}, \mathcal{M}, \mathcal{O})$ in a language-agnostic
schema: a lexicon $\mathcal{L}$ mapping surface stems to semantic lemmas
(Gilbertese \emph{noori-}\,$\leftrightarrow$\,`see'); morphology and
agreement rules $\mathcal{M}$ (object-number verb suffixes \emph{-a}/\emph{-i};
Bantu noun-class concord prefixes \emph{u-}/\emph{ba-}); and ordering
constraints $\mathcal{O}$ (verb--object--subject in Gilbertese). Nothing in
the schema names a language.

\paragraph{Certification.} A frontier LLM (Claude 5 Opus) proposes an
initial hypothesis space---candidate morpheme segmentations, lexical
mappings, and ordering constraints---expressed as an Answer Set Program.
The proposal step requires no reliability: given this space, the solver
(\texttt{clingo}; \citealp{gebser2014}) searches for a minimal rule set
that (i) regenerates every context translation pair exactly and (ii)
uniquely derives the gold translation of the target question. Any rule set
violating a single context pair is rejected. A rule set satisfying both
constraints is \emph{solver-certified}. Certification does not claim the
recovered grammar is the unique linguistic analysis of the language; it
guarantees that, within the generated space, the rules account for all
observed data and the answer. An offline induction mode enumerates all
consistent rule sets, flagging puzzles whose context underdetermines the
grammar. When certification fails, the solver returns the violated context
pair as a counterexample, which is fed back to the proposer in a bounded
repair loop. \TODO{report the loop quantitatively in
\S\ref{sec:results-generation}}

\paragraph{From rules to interventions.} Each certified rule $R$ is
compiled into a counterfactual minimal pair: a base prompt
$\mathbf{x}_{\text{base}}$ and a source prompt $\mathbf{x}_{\text{source}}$
differing only in the feature $R$ governs (e.g., singular vs.\ plural
object), with both sides re-certified. Across the four questions this
yields $N{=}29$ rules: 6 (Gilbertese Q3), 7 (Ndebele Q1), and 8 each
(Tshiluba Q7/Q8); the full schema, rule lists, and pairs are in
Appendix~\ref{app:rules}.

\section{Causal Validation}
\label{sec:validation}

All validation machinery is standard; we specify our configuration and one
computational shortcut.

\paragraph{Activation patching.} For rule $R$ with pair
$(\mathbf{x}_{\text{base}}, \mathbf{x}_{\text{source}})$, we sweep residual-stream
activations over every layer $\ell$ and token position $p$, patch
$\mathbf{h}^{(\ell,p)}_{\text{base}} \leftarrow
\mathbf{h}^{(\ell,p)}_{\text{source}}$, and report the normalized causal
indirect effect on the logit difference $\Delta$ at the first divergent
answer token,
$\mathrm{CIE}_{\mathrm{norm}}(\ell,p) =
(\Delta_{\text{patched}} - \Delta_{\text{base}}) /
(\Delta_{\text{source}} - \Delta_{\text{base}})$,
where $1$ is full causal recovery; the normalization is not capped, so
overshoot can exceed $1$ \citep{vig2020,meng2022,zhang2024}.

\paragraph{Subspace interventions.} To localize each rule below the level
of a full residual-stream patch, we use interchange interventions on a
rank-$k$ orthonormal subspace ($k \in \{1,2,4,8\}$): the standard low-rank
DAS parameterization \citep{geiger2024das,wu2023boundless}, as implemented
in pyvene \citep{wu2024pyvene} and used by ReFT \citep{wu2024reft} and
AxBench \citep{wu2025axbench}. We claim no methodological novelty here.
Our one deviation is computational: instead of optimizing the interchange
objective through the frozen network, we fit the basis $W$ against a local
linear logit surrogate (ridge regression on counterfactual activation
pairs) and then validate the resulting subspace by causal intervention on
the original model (Appendix~\ref{app:subspace}). \TODO{report agreement
between surrogate-fitted and interchange-trained subspaces at matched
sites and ranks; if agreement is high, the surrogate is a cheap
approximation, otherwise we fall back to standard training}

\paragraph{Circuits and composition.} Path patching
\citep{wang2022ioi,goldowskydill2023} traces which attention heads
transmit each rule's subspace signal, retaining edges with normalized
effect $\ge 0.15$. For composition, each pair of independent rules
$(A,B)$ is tested with a four-prompt bundle (base, source $A$, source $B$,
joint $A{+}B$): we apply interventions in both orders and measure
order-sensitivity by KL divergence between output distributions, plus
recovery of the joint intervention (Appendix~\ref{app:commutativity}).

\section{Experimental Setup}
\label{sec:setup}

We evaluate four questions from three UKLO puzzles: 2023 Gilbertese Q3,
2019 Ndebele Q1, and Tshiluba Q7/Q8,\footnote{Puzzle sources are linked in
Appendix~\ref{app:rules}; we study only questions \texttt{Llama-3.1-8B}
answers exactly correctly in context, so interventions probe a mechanism
that succeeds (see Limitations).} spanning Austronesian VOS morphology and
Bantu SVO concord systems. All experiments use
\texttt{meta-llama/Llama-3.1-8B} \citep{grattafiori2024} via NNsight/NDIF
\citep{fiottokaufman2025}: 238{,}975 patched forward evaluations, subspace
analyses at 145 causal sites, 5{,}069 path-patching traces, and 31
paired-rule intervention runs (16 distinct rule pairs). Gilbertese
supports $N{=}18$ counterfactual training pairs (ranks up to 8); Ndebele
and Tshiluba support $N{=}2$ (ranks capped at 2).

\section{Results}
\label{sec:results}

\subsection{Hypothesis generation}
\label{sec:results-generation}

\TODO{This subsection reports the generation-side experiments now in
progress; placeholders below state what each number is. (a) Certification
rate and repair-iteration counts per proposer model (Claude 5 Opus vs.\ an
open-weight proposer), with a taxonomy of solver counterexamples. (b)
Uniqueness: for each question, whether the certified rule set is unique
among all consistent rule sets (solver induction mode) or one of $m$
alternatives; underdetermined items flagged. (c) Hypothesis-space
ablation: a generic schema-only space vs.\ the LLM-proposed space ---
does proposer quality change which rules certify?}

\subsection{Causal validation}
\label{sec:results-validation}

\begin{table}[t]
\centering
\small
\begin{tabular}{@{}llrl@{}}
\toprule
Question & Rule & $|\mathrm{CIE}|$ & Site \\
\midrule
Gilb.\ Q3 & \texttt{tense}       & 3.39  & L00, P001 \\
          & \texttt{obj\_lemma}  & 0.98  & L30, P195 \\
Ndeb.\ Q1 & \texttt{yes}         & 1.20  & L03, P003 \\
          & \texttt{verb}        & 1.00  & L29, P299 \\
Tshi.\ Q7 & \texttt{tense}       & 24.00 & L00, P253 \\
          & \texttt{verb}        & 1.04  & L20, P263 \\
Tshi.\ Q8 & \texttt{obj\_small}  & 2.54  & L01, P000 \\
          & \texttt{subj\_lemma} & 0.99  & L31, P261 \\
\bottomrule
\end{tabular}
\caption{Peak causal sites for representative rules (full sweep in
Appendix~\ref{app:sweeps}). In each question, morphosyntactic rules (first
row) peak in early layers; lemma selection (second row) peaks near the
output.}
\label{tab:localization}
\end{table}

\begin{figure}[t]
\centering
% Drop the layer x position CIE heatmap panel here (early morphosyntactic
% rule vs. late lemma rule), e.g. from the workshop version's Fig. 10:
% \includegraphics[width=\columnwidth]{figures/heatmaps.pdf}
\fbox{\parbox{0.93\columnwidth}{\centering\vspace{1.0cm}
\TODO{port layer$\times$position $\mathrm{CIE}$ heatmaps (early
morphosyntax vs.\ late lemma) from the workshop version}
\vspace{1.0cm}}}
\caption{Layer$\times$position causal maps for two certified rules of the
same puzzle: morphosyntax localizes early, lemma selection late.}
\label{fig:heatmaps}
\end{figure}

\paragraph{Depth stratification.} Across all 238{,}975 grid evaluations,
every certified rule exhibits localized causal hotspots, with a consistent
pattern across the three languages (Table~\ref{tab:localization},
Figure~\ref{fig:heatmaps}): tense, agreement, and polarity markers peak in
layers L0--L4 at morpheme positions, while noun and verb lemma selection
consolidates in output-adjacent layers (L20--L31).

\paragraph{Decodable $\neq$ used.} Linear probes at the final prompt token
recover every certified rule almost perfectly ($\approx$100\% best-layer
accuracy, with selectivity controls), yet one-dimensional patches at that
token yield near-zero causal effect ($|\mathrm{CIE}|<0.05$;
Appendix~\ref{app:lasttoken}). Rule information is computed and used at
earlier morpheme positions---a caution for pipelines that read knowledge
off the last token.

\paragraph{Low-rank causal control.} Subspace interventions concentrate
each rule's effect in a compact basis: rank 8 recovers
$\mathrm{CIE}=0.575$ for Gilbertese \texttt{tense}; rank 2 reaches
$0.50$ (Ndebele \texttt{part}) and $0.76$ (Tshiluba \texttt{subj\_num});
remaining rules show stable effects in $[0.29,0.45]$ at ranks $\le 2$
(Appendix~\ref{app:subspace}). Pruned circuits from path patching retain
mean $|\mathrm{CIE}|$ of 0.77--1.04 while removing over 85\% of candidate
edges (Appendix~\ref{app:circuits}).

\paragraph{Independent rules compose.} Sequential interventions on
independent rule pairs commute on 31/31 pair--site evaluations (16 rule
pairs; max order-KL $\le 0.0118$ nats) and recover the corresponding joint
interventions (Appendix~\ref{app:commutativity}). Controls: linear
subspaces match or beat 2-layer MLP surrogates at 53.5--79.2\% of sites;
random bases have no effect ($|\mathrm{CIE}|<0.01$); commutativity holds
over the top-50 output tokens ($D_{\mathrm{KL}} \le 0.024$ nats), not just
the answer logits.

\section{Conclusion}
\label{sec:conclusion}

We showed that the hypothesis-generation step of causal interpretability
can be delegated to task structure: a solver derives and certifies the
grammatical rules a Rosetta puzzle requires, and standard interventions
then test each certified rule inside the model. On
\texttt{Llama-3.1-8B}, certified rules are implemented in depth-stratified,
low-rank linear subspaces that compose additively---and are linearly
decodable at the final token without being causally used there. Wherever a
domain admits an executable hypothesis format and a checker, the same
certify-then-validate loop turns ``the model seems to know $X$'' into two
auditable claims: $X$ provably accounts for the task data, and the model
causally computes it.

\section*{Limitations}
\label{sec:limitations}

Our study is deep rather than broad: four questions, three languages, one
8B model. We analyze only questions the model answers exactly correctly,
which is necessary to study a succeeding mechanism but conditions every
finding on success; mechanisms behind near-misses may differ. Ndebele and
Tshiluba rules admit only $N{=}2$ counterfactual pairs, capping subspace
rank at 2 and limiting statistical power; results there are directional.
The subspace bases are fitted against a local linear surrogate rather than
by interchange training through the network; although every subspace is
subsequently validated by causal intervention on the model, agreement with
standard end-to-end training remains to be quantified
\TODO{resolve with the planned agreement experiment}. Single-token answer
evaluation sidesteps multi-token generation. The grammar schema covers
concatenative morphology and consistent ordering; fusional morphology,
reduplication, and phonological alternation are outside its current scope.
Finally, evidence for additive composition is consistency across 16 rule
pairs, not a proof; larger rule inventories could reveal interference.

\section*{Ethics Statement}
\label{sec:ethics}

This work analyzes the internal mechanisms of a publicly released language
model on educational linguistics puzzles; it introduces no new model or
generation capability and no data about individuals. UKLO materials are
used under their free-for-educational-use terms, with puzzle sources
credited. The principal societal risk we foresee is epistemic:
knowledge-extraction pipelines can lend unearned authority to spurious
``discoveries'' read out of models. Our framework targets exactly this
failure mode---every extracted claim must be solver-certified against data
and causally validated in the model---and we report the failure of the
cheaper last-token recipe so practitioners do not over-trust it.


\bibliography{custom}

\appendix
% Appendices (not counted toward the short-paper page limit).
% Each section marks the material to port verbatim from the workshop
% version's LaTeX source (the NeurIPS Interp4Discovery submission).

\section{Grammar Schema, Certified Rules, and Minimal Pairs}
\label{app:rules}

\TODO{Port from the workshop version: formal schema
$G = (\mathcal{L}, \mathcal{M}, \mathcal{O})$ and the certification
formulation (its App.~B); the full certified rule lists for all four
questions ($N{=}29$: 6 Gilbertese Q3, 7 Ndebele Q1, 8 each Tshiluba
Q7/Q8); puzzle sources and full prompts (its App.~E); and the
counterfactual minimal-pair tables (its App.~L). Note: fix the stale
``$N=27$'' headline from the workshop version's App.~B.3 --- the correct
count is 29.}

\section{Subspace Fitting Details}
\label{app:subspace}

\TODO{Port the surrogate objective and training details (its App.~C),
restated in standard low-rank DAS terminology, plus the rank-validation
table (its App.~G, Table~2). Add the planned agreement experiment:
surrogate-fitted vs.\ interchange-trained bases at matched sites and
ranks.}

\section{Full Patching Sweeps and Last-Token Comparison}
\label{app:sweeps}
\label{app:lasttoken}

\TODO{Port the per-rule peak-CIE table (its App.~F, Table~1), the
layer$\times$position heatmaps (its App.~K, Fig.~10), and the 1D
last-token probing vs.\ sweep comparison (its App.~J, Table~6).}

\section{Circuits}
\label{app:circuits}

\TODO{Port the pruned circuit graphs and head-level effects (its App.~K,
Fig.~11) and the MDL/faithfulness analysis (Fig.~12) if space warrants.}

\section{Commutativity Details}
\label{app:commutativity}

\TODO{Port the commutativity and KL formulations (its App.~D) and the
summary and pairwise tables (its App.~I, Tables~4--5), plus the ablation
summary (its App.~H, Table~3).}


\end{document}
